\documentclass[aspectratio=169,10pt]{beamer}

\usefonttheme{professionalfonts}
\setbeamertemplate{navigation symbols}{}

\usepackage{graphicx}
\usepackage{booktabs}
\usepackage{hyperref}
\usepackage{amsmath}
\usepackage{bm}
\usepackage{ragged2e}
\usepackage{xcolor}
\usepackage{array}
\usepackage{subcaption}

\definecolor{SCUGreen}{RGB}{8,79,52}
\definecolor{SCULight}{RGB}{223,233,227}
\definecolor{SCUDot}{RGB}{35,137,79}
\definecolor{SCUText}{RGB}{41,63,70}

\setbeamercolor{background canvas}{bg=white}
\setbeamercolor{normal text}{fg=SCUText,bg=white}
\setbeamercolor{title}{fg=SCUGreen}
\setbeamercolor{subtitle}{fg=SCUGreen}
\setbeamercolor{frametitle}{fg=SCUGreen,bg=SCULight}
\setbeamercolor{itemize item}{fg=SCUDot}
\setbeamercolor{itemize subitem}{fg=SCUDot}
\setbeamercolor{enumerate item}{fg=SCUDot}

\setbeamerfont{title}{size=\Large,series=\bfseries}
\setbeamerfont{subtitle}{size=\large}
\setbeamerfont{frametitle}{size=\large,series=\bfseries}
\setbeamerfont{normal text}{size=\normalsize}

\setbeamertemplate{footline}{}
\setbeamertemplate{headline}{}
\setbeamersize{text margin left=1.15cm,text margin right=1.15cm}
\setlength{\leftmargini}{1.2em}
\setlength{\leftmarginii}{1.15em}
\setbeamertemplate{itemize item}{\color{SCUDot}\small$\bullet$}
\setbeamertemplate{itemize subitem}{\color{SCUDot}\scriptsize$\bullet$}
\setbeamertemplate{itemize subsubitem}{\color{SCUDot}\tiny$\bullet$}
\setbeamertemplate{frametitle}{
  \nointerlineskip
  \begin{beamercolorbox}[wd=\paperwidth,ht=3.1ex,dp=1.1ex,leftskip=1.1em]{frametitle}
    \usebeamerfont{frametitle}\insertframetitle
  \end{beamercolorbox}
}

\newcommand{\pipelinebox}[1]{
  \begin{beamercolorbox}[wd=\linewidth,sep=0.8em,rounded=true,center]{frametitle}
    \color{SCUGreen}\bfseries #1
  \end{beamercolorbox}
}

\newcommand{\sectiondivider}[2]{
  \begin{frame}[plain]
    \vspace{1.75cm}
    {\color{SCUGreen}\bfseries\fontsize{15.2}{18}\selectfont #1\par}
    \vspace{0.42cm}
    {\color{SCUGreen}\fontsize{13.2}{15.6}\selectfont #2\par}
    \vspace{0.28cm}
    {\color{SCUGreen}\rule{0.92\linewidth}{0.55pt}\par}
    \vfill
  \end{frame}
}

\newcommand{\imageorplaceholder}[2]{
  \IfFileExists{#1}{
    \includegraphics[width=\linewidth,height=#2,keepaspectratio]{#1}
  }{
    \fbox{
      \parbox[c][#2][c]{0.92\linewidth}{
        \centering
        \textbf{Placeholder}\\[0.4em]
        \path{#1}
      }
    }
  }
}

\title{From Fourier Series to Spherical Harmonics}
\subtitle{View-Dependent Appearance in 3D Gaussian Splatting $\left( 2026/06/03 \right)$}
\author{Heyang (Stefan) Zhang}
\institute{Sichuan University - Pittsburgh Institute}
\date{}

\begin{document}

\begin{frame}[plain]
  \vspace{1.55cm}
  {\color{SCUGreen}\bfseries\fontsize{15.2}{18}\selectfont \inserttitle\par}
  \vspace{0.42cm}
  {\color{SCUGreen}\fontsize{13.2}{15.6}\selectfont \insertsubtitle\par}
  \vspace{0.28cm}
  {\color{SCUGreen}\rule{0.96\linewidth}{0.55pt}\par}
  \vfill
  {\centering\fontsize{11.8}{14}\selectfont \insertauthor\par}
  \vspace{0.75cm}
  {\centering\fontsize{11}{13}\selectfont \insertinstitute\par}
  \vspace{1.15cm}
\end{frame}

\begin{frame}[t]{Contents}
  \large
  \vspace{1.5cm}
  \begin{itemize}
    \item Section I: Background and Fundamental Theory
    \item Section II:  Spherical Harmonics
    \item Section III: Directional Color as a Visual Signal
    \item Section IV: Application in 3D Gaussian Splatting
    \item Section V: Demo and Evaluation
    \item Conclusion
  \end{itemize}
\end{frame}

\section{Background}
\sectiondivider{Section I: Background and Fundamental Theory}{From ordinary signals to basis-function expansions}

\begin{frame}[t]{Background}
  \justifying
\begin{itemize}
  \item Fourier analysis originated from the study of heat transfer in the early 19th century.
  \item Joseph Fourier proposed that complex periodic signals could be represented by sums of simple sine and cosine waves.
  \item This idea later became a foundation of signal processing, communication systems, image analysis, and modern visual computing.
  \item The core insight is simple: complicated signals can be decomposed into simpler basis functions.
\end{itemize}

  \vspace{6.75em}
  \pipelinebox{Complex Signal $\rightarrow$ Sine and Cosine Bases $\rightarrow$ Frequency Coefficients}
\end{frame}

\begin{frame}{Fourier Series}
A periodic signal can be represented as a weighted sum of sine and cosine functions:

\[
f(x)=a_0+\sum_{n=1}^{\infty}
\left(a_n\cos(nx)+b_n\sin(nx)\right)
\]

where:
\begin{itemize}
  \item $a_0$ represents the average level of the signal.
  \item $a_n$ and $b_n$ are Fourier coefficients.
  \item $\cos(nx)$ and $\sin(nx)$ are basis functions.
  \item $n$ controls the frequency of each component.
\end{itemize}

\end{frame}

\begin{frame}[t]{Approximating a Signal by Increasing $n$}
  \justifying
  \begin{columns}[T]
    \column{0.25\linewidth}
\begin{itemize}
  \item Using only low-frequency terms gives a rough and smooth approximation.
  \item Adding more terms allows the representation to capture sharper changes and finer details.
\end{itemize}
    \column{0.75\linewidth}
    \begin{figure}
    \centering
    \imageorplaceholder{figures/fourier_square_wave_approx.png}{0.86\textheight}
    \caption{Square-wave approximation with increasing Fourier terms}
    \end{figure}
  \end{columns}
\end{frame}

\begin{frame}[t]{Fourier Ideas in 2D Images}
  \justifying
  \begin{columns}[T]
    \column{0.5\linewidth}
\begin{itemize}
  \item For object boundaries, a closed contour can also be represented in polar form:
  \[
  r=r(\theta)
  \]
  \item Then Fourier series can approximate this boundary by adding more frequency components.
\end{itemize}

    \column{0.5\linewidth}
    \begin{figure}
    \centering
    \imageorplaceholder{figures/polar_dumbbell_r_costheta.png}{0.7\textheight}
    \caption{A dumbbell-shape in polar coordinates}
    \end{figure}
  \end{columns}
\end{frame}

% =========================
% Section I - Extra Page
% =========================
\begin{frame}{Two-Variable Fourier Series}
A two-variable signal can be written as:
\[
z=f(x,y)
\]

where $x$ and $y$ are spatial coordinates, and $z$ is the signal value or surface height.

For a periodic 2D signal, the Fourier series can be written as:

\[
f(x,y)
=
\sum_{m=-M}^{M}
\sum_{n=-N}^{N}
c_{mn}
e^{j(mx+ny)}
\]

Equivalently, it can be expressed using sine and cosine terms:

\[
\begin{aligned}
f(x,y)\approx
\sum_{m=0}^{M}\sum_{n=0}^{N}
\big[
 A_{mn}\cos(mx)\cos(ny) 
+ B_{mn}\cos(mx)\sin(ny) \\
+ C_{mn}\sin(mx)\cos(ny) 
+ D_{mn}\sin(mx)\sin(ny)
\big]
\end{aligned}
\]
\end{frame}

\begin{frame}{Basis Functions of a Two-Variable Fourier Series}
\begin{figure}
\centering
\begin{subfigure}{0.23\textwidth}
\centering
\imageorplaceholder{figures/two_variable_basis_n1.png}{0.53\textheight}
\caption{$\cos(x)\cos(y)$}
\end{subfigure}
\hfill
\begin{subfigure}{0.23\textwidth}
\centering
\imageorplaceholder{figures/two_variable_basis_n2.png}{0.53\textheight}
\caption{$\cos(x)\cos(2y)$}
\end{subfigure}
\hfill
\begin{subfigure}{0.23\textwidth}
\centering
\imageorplaceholder{figures/two_variable_basis_n3.png}{0.53\textheight}
\caption{$\cos(x)\cos(3y)$}
\end{subfigure}
\hfill
\begin{subfigure}{0.23\textwidth}
\centering
\imageorplaceholder{figures/two_variable_basis_n4.png}{0.53\textheight}
\caption{$\cos(x)\cos(4y)$}
\end{subfigure}
\caption{Basis surfaces with increasing frequency in the $y$ direction}
\end{figure}
\end{frame}

 \begin{frame}{Two-Variable Fourier Series for Surface Signals}
\begin{columns}
\column{0.5\textwidth}
\begin{center}
\begin{figure}
\centering
\imageorplaceholder{figures/two_variable_surface.png}{1\textheight}
\caption{Two-variable surface signal}
\end{figure}
\end{center}
\column{0.5\textwidth}
\begin{center}
\begin{figure}
\centering
\imageorplaceholder{figures/two_variable_fit_grid.png}{1\textheight}
\caption{Increasing 2D Fourier terms}
\end{figure}
\end{center}
\end{columns}
 \end{frame}



\section{Spherical Harmonics}
\sectiondivider{Section II: Spherical Harmonics}{Extending Fourier representation from a plane to a sphere}



\begin{frame}{Fourier-Like Expansion in Spherical Coordinates}
  \justifying
\begin{columns}
\column{0.5\textwidth}
In spherical coordinates, a 3D surface can be described by a radius function:
\[
r=r(\theta,\phi)
\]
where:
\begin{itemize}
  \item $r$ is the distance from the origin.
  \item $\theta$ is the polar angle.
  \item $\phi$ is the azimuth angle.
\end{itemize}
The corresponding Cartesian coordinates are:
\[
x=r(\theta,\phi)\sin\theta\cos\phi
\]
\[
y=r(\theta,\phi)\sin\theta\sin\phi
\]
\[
z=r(\theta,\phi)\cos\theta
\]
    \column{0.4\linewidth}
    \begin{figure}
    \centering
    \imageorplaceholder{figures/spherical_coordinates.png}{0.68\textheight}
    \caption{Spherical coordinates with radius, polar angle, and azimuth angle}
    \end{figure}
  \end{columns}
  \end{frame}



\begin{frame}{Fourier-Like Expansion in Spherical Coordinates}
\justifying
    
Therefore, the basis expansion changes its domain from a plane to a sphere.

For a 2D signal on a plane:
\[
f(x,y)\approx
\sum_{m=-M}^{M}
\sum_{n=-N}^{N}
c_{mn}\Phi_{mn}(x,y)
\]


For a directional signal on a sphere:

\[
r(\theta,\phi)
=
\sum_{l=0}^{L}\sum_{m=-l}^{l}
c_l^m Y_l^m(\theta,\phi)
\]

where:
\begin{itemize}
  \item $l$ is the degree, controlling the frequency level.
  \item $m$ is the order, controlling angular variation.
  \item $Y_l^m(\theta,\phi)$ is a spherical harmonic basis function.
  \item $c_l^m$ is the spherical harmonic coefficient.
\end{itemize}

\end{frame}

\begin{frame}{Spherical Harmonics and SH Coefficients}
\begin{columns}
\column{0.5\textwidth}
Spherical harmonics are basis functions defined on the sphere:
\[
Y_l^m(\theta,\phi)
\]
\small
\[
Y_l^m(\theta,\varphi)=
\begin{cases}
\sqrt{2}\,K_l^m\cos(m\varphi)\,P_l^m(\cos\theta),
& m>0\\[0.6em]
\sqrt{2}\,K_l^m\sin(-m\varphi)\,P_l^{-m}(\cos\theta),
& m<0\\[0.6em]
K_l^0P_l^0(\cos\theta),
& m=0
\end{cases}
\]
\[
P_n(x)=
\frac{1}{2^n n!}
\frac{d^n}{dx^n}
\left[
(x^2-1)^n
\right]
\]
\[
P_l^m(x)=
(-1)^m(1-x^2)^{m/2}
\frac{d^m}{dx^m}
\left(
P_l(x)
\right)
\]

\[
K_l^m=
\sqrt{
\frac{(2l+1)}{4\pi}
\frac{(l-|m|)!}{(l+|m|)!}
}
\]
\column{0.5\textwidth}

\begin{center}
\begin{figure}
\centering
\imageorplaceholder{figures/sh_basis_grid.png}{0.58\textheight}
\caption{SH basis functions arranged by degree $L$}
\end{figure}
\end{center}
\end{columns}
\end{frame}


\begin{frame}{Spherical Harmonics and SH Coefficients}

The coefficient measures how much each basis function contributes:

\[
c_l^m
=
\int_{0}^{2\pi}
\int_{0}^{\pi}
f(\theta,\phi)
\overline{Y_l^m(\theta,\phi)}
\sin\theta
\,d\theta\,d\phi
\]
\vspace{2.2em}
\begin{center}
\begin{figure}
\centering
\imageorplaceholder{figures/sh_coefficient_projection.png}{0.3\textheight}
\caption{A signal reconstructed as weighted SH basis functions}
\end{figure}
\end{center}
\end{frame}

\begin{frame}{Increasing SH Degree Improves Approximation}
\begin{center}
\begin{figure}
\centering
\imageorplaceholder{figures/sh_approximation_grid.png}{0.8\textheight}
\caption{Higher SH degree captures finer angular detail}
\end{figure}
\end{center}
\end{frame}



\section{Directional Color as a Visual Signal}
\sectiondivider{Section III: Directional Color as a Visual Signal}{From color values to signal representation}

\begin{frame}{Color Can Be Represented as a Signal}
\begin{columns}
\column{0.52\textwidth}

\begin{itemize}
  \item Color is not only a visual property. It can also be treated as a measurable signal.
  \item For a fixed image point, the color value can be written as:
  \[
  C = (R,G,B)
  \]
  \item If color changes with position, it becomes a spatial signal:
  \[
  C(x,y) = \big(R(x,y),G(x,y),B(x,y)\big)
  \]
  \item If color changes with viewing direction, it becomes a directional signal:
  \[
  C(\theta,\phi)
  \]
\end{itemize}


\column{0.43\textwidth}

\begin{figure}
\centering
\imageorplaceholder{figures/color_as_signal_grayscale.png}{0.52\textheight}
\caption{Gray values represented as a signal}
\end{figure}

\end{columns}
\end{frame}

\begin{frame}{From Directional Distance to Color Signal}

\begin{columns}
\column{0.45\textwidth}
\begin{center}
\begin{figure}
\centering
\imageorplaceholder{figures/directional_distance_signal_custom.png}{0.9\textheight}
\caption{Directional distance signal}
\end{figure}
\end{center}
\column{0.1\textwidth}
\begin{center}
\hspace{0.5cm}
$\Longrightarrow$
\hspace{0.5cm}
\end{center}
\column{0.45\textwidth}
\begin{center}
\begin{figure}
\centering
\imageorplaceholder{figures/color_mapped_directional_signal.png}{0.9\textheight}
\caption{Color-mapped signal on the sphere}
\end{figure}
\end{center}
\end{columns}
\end{frame}

\begin{frame}{Multi-Channel Signals and Directional Anisotropy}
\begin{columns}
\column{0.40\textwidth}
\begin{itemize}
  \item Color can be represented by three directional channels:
  \[
  C(\theta,\phi)
  =
  \big(
  R(\theta,\phi),
  G(\theta,\phi),
  B(\theta,\phi)
  \big)
  \]
  \item Brightness, transparency, and other visual attributes can also be modeled as signals:
  \[
  I(\theta,\phi),\quad
  \alpha(\theta,\phi)
  \]
  \item If a property changes with direction, it is anisotropic.
\end{itemize}

\column{0.6\textwidth}
\begin{figure}
\centering
\imageorplaceholder{figures/multichannel_directional_signal_custom.png}{1\textheight}
\caption{\scriptsize RGB, brightness, and opacity as directional signals}
\end{figure}

\end{columns}
\end{frame}



\section{Application in 3D Gaussian Splatting}
\sectiondivider{Section IV: Application in 3D Gaussian Splatting}{Using spherical harmonics to model view-dependent appearance}

\begin{frame}{SH Color Representation}
\begin{columns}
\column{0.4\textwidth}
Spherical harmonics are used to represent directional color:
\[
\mathbf{C}_i(\bm d)
=
\sum_{l=0}^{L}
\sum_{m=-l}^{l}
\mathbf{c}_{i,l}^{m}
Y_l^m(\bm d)
\]
where:
\begin{itemize}
  \item $\bm d$ is the viewing direction from the Gaussian to the camera.
  \item $Y_l^m(\bm d)$ is the SH basis value.
  \item $\mathbf{c}_{i,l}^{m}$ is the learned RGB coefficient.
\end{itemize}
For RGB color:
\[
\mathbf{c}_{i,l}^{m}
=
\left(
c_{R,i,l}^{m},
c_{G,i,l}^{m},
c_{B,i,l}^{m}
\right)
\]
\column{0.6\textwidth}
\begin{figure}
\centering
\imageorplaceholder{figures/sh_color_six_views.png}{0.95\textheight}
\caption{View-dependent color under SH coefficients}
\end{figure}

\end{columns}
\end{frame}

\begin{frame}{3D Gaussian Primitive}
\begin{columns}
\column{0.4\textwidth}
\[
\mathcal{G}_i =
\{\bm{\mu}_i,\Sigma_i,\alpha_i,\mathbf{c}_i\}
\]
where:
\begin{itemize}
  \item $\bm{\mu}_i$ is the 3D position.
  \item $\Sigma_i$ controls anisotropic shape.
  \item $\alpha_i$ is opacity.
  \item $\mathbf{c}_i$ represents color information.
\end{itemize}
\column{0.6\textwidth}
\begin{figure}
\centering
\imageorplaceholder{figures/gaussian_ellipsoid_six_views.png}{0.95\textheight}
\caption{An anisotropic 3D Gaussian primitive}
\end{figure}
\end{columns}
\end{frame}

\begin{frame}{From One Gaussian to a 3D Scene}
A single Gaussian only represents a small local primitive.
\vspace{0.2cm}
By combining thousands or even millions of Gaussians, a complex object or scene can be approximated:
\[
\mathcal{S}
=
\{\mathcal{G}_1,\mathcal{G}_2,\ldots,\mathcal{G}_N\}
\]
where each Gaussian has its own:
\[
\mathcal{G}_i =
\{\bm{\mu}_i,\Sigma_i,\alpha_i,\mathbf{c}_i\}
\]
During rendering, these Gaussians are projected to the image plane and blended together:
\[
\mathbf{I}(u,v)
\approx
\sum_{i=1}^{N}
T_i(u,v)\,\alpha_i'(u,v)\,\mathbf{C}_i(\bm d)
\]
where $T_i$ represents accumulated transmittance from previously blended Gaussians.
\end{frame}

\begin{frame}{The Pipeline of 3D Gaussian Splatting}
\begin{center}
\begin{figure}
\centering
\imageorplaceholder{figures/3dgs_pipeline_fig2_original.png}{0.82\textheight}
\caption{Original optimization pipeline of 3D Gaussian Splatting}
\end{figure}
\end{center}
\end{frame}

\begin{frame}{Examples of 3DGS}
\begin{columns}
\column{0.45\textwidth}
\begin{center}
\begin{figure}
\centering
\imageorplaceholder{figures/000050_overlay.png}{0.8\textheight}
\caption{Novel view with directional color}
\end{figure}
\end{center}

\column{0.45\textwidth}
\begin{center}
\begin{figure}
\centering
\imageorplaceholder{figures/000022_overlay.png}{0.8\textheight}
\caption{Novel view with directional color}
\end{figure}
\end{center}
\end{columns}
\end{frame}

\section{Demo and Evaluation}
\sectiondivider{Section V: Demo and Evaluation}{A panorama-based demonstration of spherical harmonic reconstruction}

\begin{frame}[t]{Demo Design}
  \justifying
  \begin{columns}[T]
    \column{0.52\linewidth}
    \begin{itemize}
      \item The user selects a 360$^\circ$ panorama image from a Windows desktop GUI.
      \item The application computes spherical harmonic coefficients for the RGB channels.
      \item This demo is a simplified visualization of the directional-color idea used in 3DGS.
      \item The demo displays the original panorama and reconstructed results under different SH orders.
    \end{itemize}

    \column{0.42\linewidth}
    \pipelinebox{
      Select 360$^\circ$ Image\\[0.2em]
      $\Downarrow$\\[0.2em]
      Compute SH Coefficients\\[0.2em]
      $\Downarrow$\\[0.2em]
      Compare Reconstructions
    }
  \end{columns}
\end{frame}




\begin{frame}[t]{Conclusion}
  \justifying
  \begin{itemize}
    \item Fourier series show how complex signals can be represented by weighted basis functions.
    \item Spherical harmonics extend this idea from one-dimensional or planar signals to directional signals on a sphere.
    \item In 3D Gaussian Splatting, SH coefficients are used to model view-dependent RGB color.
    \item The panorama demo illustrates how increasing the SH order improves directional signal reconstruction.
  \end{itemize}
\end{frame}

\begin{frame}[t]{Team Members and Responsibilities}
\centering
\begin{tabular}{l l}
\toprule
\textbf{Name} & \textbf{Role / Contribution} \\
\midrule
Heyang Zhang & Presentation, Code Integration \\
Dexi Liu & Research on 3D Gaussian Splatting, Data Preparation \\
Shuyu Wang & Demo Implementation, Visualization \\
Hongwei Zeng & Analysis of Directional Color Signals, Slide Preparation \\
Rui Wang  & Supporting Research, Literature Review, Figures Preparation \\
\bottomrule
\end{tabular}
\end{frame}

\end{document}
